% Separate study blocks retain their own reference values. They must not be
% pooled into a single CAST row. Sources: tab:constraint, tab:exp1, and the
% reward-field block of tab:locus in the ICLR draft. Older provisional hinge
% and projection rows are not substituted for the new three-seed comparison.
\begin{table}[t]
\centering
\caption{Constraint and reward-field ablations. (a) Reported three-seed
five-task LIBERO factorial; mean base success is .673. (b) The constraint
comparison uses three training seeds, 300-episode square evaluations, and powered
LIBERO evaluations. (c) Reward-field substitutions, with square seed ranges.
Each study block retains its own reported CAST reference; values are not
pooled across blocks. Mug and ketchup denote tasks 65 and 32.
Bold marks the best result per column within each study block.}
\label{tab:ablations}
\footnotesize
\setlength{\tabcolsep}{4pt}
\begin{tabular}{@{}llc@{}}
\toprule
\multicolumn{3}{@{}l}{\emph{(a) Step clipping versus base restoration}} \\
Step clip & Restoration & Five-task mean \\
\midrule
Off & Off & .000 \\
On & Off & .000 \\
Off & On & .657 \\
On & On (full CAST) & \textbf{.703} \\
\bottomrule
\end{tabular}
\par\smallskip
\begin{tabular}{@{}lccc@{}}
\toprule
\multicolumn{4}{@{}l}{\emph{(b) Alternative constraint implementations}} \\
Constraint & Square & Mug & Ketchup \\
\midrule
Loss-space hinge & .28 & .77 & .65 \\
Hard target projection & \textbf{.96} & \textbf{.93} & .54 \\
Full CAST & .94 & .78 & \textbf{.68} \\
\midrule
\multicolumn{4}{@{}l}{\emph{(c) Reward field, with the constraint fixed}} \\
Field & Square & Mug & Ketchup \\
\midrule
Gradient & \textbf{.912--.930} & .600 & \textbf{.627} \\
Top-$k$ & .525--.562 & .577 & .341 \\
Value-tilted & .560--.584 & \textbf{.697} & .378 \\
\bottomrule
\end{tabular}
\end{table}
